\section*{Chapter \ref{ch:g11:evolution-mechanisms} --- Mechanisms of Evolution}

\begin{solution}{exo:g11:evolution-mechanisms:1}
In \omterm{def:g11:evolution-mechanisms:pool}{population} genetics, \omterm{def:g11:evolution-mechanisms:evo}{evolution} is change in \omterm{def:g11:evolution-mechanisms:evo}{allele frequencies} in a
\omterm{def:g11:evolution-mechanisms:pool}{population}'s \omterm{def:g11:evolution-mechanisms:pool}{gene pool} across generations. That definition includes selection,
drift, \omterm{def:g11:evolution-mechanisms:mutation}{mutation} and \omterm{def:g11:evolution-mechanisms:geneflow}{gene flow} --- not only the origin of new species or
dramatic morphological change.
\end{solution}

\begin{solution}{exo:g11:evolution-mechanisms:2}
\omterm{prop:g11:evolution-mechanisms:evidence}{Homologous structures} share ancestry even if functions differ (e.g.\ vertebrate
forelimbs). \omterm{prop:g11:evolution-mechanisms:evidence}{Analogous structures} share function without the same ancestral
structure, often from convergent \omterm{def:g11:evolution-mechanisms:evo}{evolution} (e.g.\ insect wing and bird wing as
flight surfaces). Phylogeny groups by homology, not mere similarity of use.
\end{solution}

\begin{solution}{exo:g11:evolution-mechanisms:3}
\omterm{def:g11:evolution-mechanisms:selection}{Natural selection} requires heritable variation in a trait, differential
survival or \omterm{def:g10:scientific-biology:properties}{reproduction} associated with that trait, and inheritance so that
offspring resemble parents. Without heritability, selection on \omterm{def:g11:mendelian-genetics:geno-pheno}{phenotypes} does
not change the next generation's genetic composition.
\end{solution}

\begin{solution}{exo:g11:evolution-mechanisms:4}
\omterm{def:g11:evolution-mechanisms:mutation}{Mutations} arise without regard to whether they will help the \omterm{def:g10:scientific-biology:organism}{organism}; most
are neutral or harmful. Adaptive \omterm{def:g11:evolution-mechanisms:evo}{evolution} still needs \omterm{def:g11:evolution-mechanisms:mutation}{mutation} as the ultimate
source of new \omterm{def:g11:mendelian-genetics:allele}{alleles} on which selection can act --- selection sorts variation,
it does not invent \omterm{def:g11:mendelian-genetics:allele}{alleles} on demand.
\end{solution}

\begin{solution}{exo:g11:evolution-mechanisms:5}
\omterm{def:g11:evolution-mechanisms:drift}{Genetic drift} changes \omterm{def:g11:evolution-mechanisms:evo}{allele frequencies} by chance sampling and is random with
respect to \omterm{def:g11:evolution-mechanisms:fitness}{fitness}; it does not systematically produce \omterm{def:g11:evolution-mechanisms:fitness}{adaptations}. \omterm{def:g11:evolution-mechanisms:selection}{Natural
selection} is non-random with respect to \omterm{def:g11:evolution-mechanisms:fitness}{fitness} differences and can produce
\omterm{def:g11:evolution-mechanisms:fitness}{adaptations}. Both change \omterm{def:g11:evolution-mechanisms:pool}{gene pools}; only selection consistently builds fit.
\end{solution}

\begin{solution}{exo:g11:evolution-mechanisms:6}
A \omterm{def:g11:evolution-mechanisms:bottle}{bottleneck} suddenly reduces \omterm{def:g11:evolution-mechanisms:pool}{population} size so survivors are a small sample
of the original \omterm{def:g11:evolution-mechanisms:pool}{gene pool}. \omterm{def:g11:evolution-mechanisms:evo}{Allele frequencies} among survivors can differ from
the original \omterm{def:g11:evolution-mechanisms:pool}{population} by pure sampling error even if every individual was
equally likely to die --- a classic drift effect, often with lost diversity.
\end{solution}

\begin{solution}{exo:g11:evolution-mechanisms:7}
If $q = 0.4$ then $p = 0.6$. Under Hardy--Weinberg, $AA$ frequency is
$p^{2} = 0.36$, $Aa$ is $2pq = 0.48$, and $aa$ is $q^{2} = 0.16$. Those are
expected \omterm{def:g11:mendelian-genetics:geno-pheno}{genotype} frequencies only if assumptions of no selection, no drift
bias, no \omterm{def:g11:evolution-mechanisms:mutation}{mutation} or migration, and random mating hold.
\end{solution}

\begin{solution}{exo:g11:evolution-mechanisms:8}
If $q^{2} = 0.01$ then $q = 0.1$ and $p = 0.9$. Carrier frequency is
$2pq = 2\times 0.9\times 0.1 = 0.18$, or $18\%$. Most copies of a rare
\omterm{def:g11:mendelian-genetics:dominance}{recessive allele} sit in heterozygotes, not in affected homozygotes.
\end{solution}

\begin{solution}{exo:g11:evolution-mechanisms:9}
Heritable variation in resistance already exists (or arises by \omterm{def:g11:evolution-mechanisms:mutation}{mutation}).
Antibiotic use kills sensitive bacteria, so resistant \omterm{def:g11:mendelian-genetics:allele}{alleles} increase in
frequency by \omterm{def:g11:evolution-mechanisms:selection}{natural selection}. Individual \omterm{def:g10:the-cell:cell}{cells} do not purposefully acquire
immunity by ``trying''; the \omterm{def:g11:evolution-mechanisms:pool}{population} evolves because survivors reproduce.
\end{solution}

\begin{solution}{exo:g11:evolution-mechanisms:10}
\omterm{def:g11:evolution-mechanisms:geneflow}{Gene flow} mixes \omterm{def:g11:mendelian-genetics:allele}{alleles} between \omterm{def:g11:evolution-mechanisms:pool}{populations} and reduces frequency differences.
Immigrants can continually reintroduce \omterm{def:g11:mendelian-genetics:allele}{alleles} that are locally less fit,
which can prevent a \omterm{def:g11:evolution-mechanisms:pool}{population} from fixing only locally optimal \omterm{def:g11:mendelian-genetics:allele}{alleles} and
can also spread adaptive \omterm{def:g11:mendelian-genetics:allele}{alleles} when they arise.
\end{solution}

\begin{solution}{pb:g11:evolution-mechanisms:1}
\textbf{1.} \omterm{def:g11:evolution-mechanisms:selection}{Natural selection} is acting; the Hardy--Weinberg assumption of no
selection is violated, so \omterm{def:g11:mendelian-genetics:geno-pheno}{genotype} frequencies need not match $p^{2}:2pq:q^{2}$
in the simple way.

\textbf{2.} With $p = 0.4$, $q = 0.6$: $BB = 0.16$, $Bb = 0.48$, $bb = 0.36$
under Hardy--Weinberg expectations.

\textbf{3.} After the change $p' = q' = 0.5$: $BB = 0.25$, $Bb = 0.50$,
$bb = 0.25$.

\textbf{4.} \omterm{def:g11:evolution-mechanisms:bottle}{Founder effect} (a form of drift): the five founders are a small
random sample of mainland \omterm{def:g11:mendelian-genetics:allele}{alleles}, so island frequencies can differ by chance.

\textbf{5.} \omterm{def:g11:evolution-mechanisms:geneflow}{Gene flow} tends to reduce genetic differences between island and
mainland by mixing \omterm{def:g11:mendelian-genetics:allele}{alleles} each generation.

\textbf{6.} \omterm{prop:g11:evolution-mechanisms:evidence}{Fossils} show long-term morphological change consistent with
cumulative generation-by-generation change; living birds show the mechanism
(\omterm{def:g11:evolution-mechanisms:evo}{allele frequency} change under selection) on short timescales.

\textbf{7.} Teleology implies intention or foresight. Better: birds with
heritable larger beaks survived and reproduced more when large seeds were
available, so \omterm{def:g11:mendelian-genetics:allele}{alleles} for larger beaks increased in frequency.
\end{solution}
